\section{Literature Review and Study of Existing Systems}

\subsection{Study of Existing Systems}
\paragraph{Manual process (current practice).}
The timetable committee prepares the timetable in a spreadsheet or word processor and circulates a
PDF by e-mail, on WhatsApp groups and on notice boards. Users read it by hand. The process costs
nothing and everyone understands it, but it supports no search, no personalisation, no clash
checking and no version control.

\paragraph{Timetable generation software: FET, UniTime and aSc Timetables.}
FET~\cite{fet} is a free, open-source timetable generator. UniTime~\cite{unitime} is an open-source,
web-based university course timetabling and student scheduling system. aSc Timetables~\cite{asc} is a
commercial school timetable generator. These tools solve the \emph{construction} problem: they
assign events to time slots and rooms under hard and soft constraints, and they export the result
as printable grids. They are used by the committee, not by students. They do not read an existing
PDF, and their published output is again a static document.

\paragraph{PDF table-extraction tools: Tabula and Camelot.}
Tabula~\cite{tabula} and Camelot~\cite{camelot} extract tables from text-based PDFs into CSV. They
handle ruled tables well, but they return raw cell strings: they do not understand time rows, lab
sub-groups, legends, merged two-hour cells or a.m./p.m.\ ambiguity, and they do not work on scanned
images without separate OCR.

\paragraph{General calendar apps: Google Calendar and Outlook.}
Calendar apps offer reminders and sharing, but every event must be entered by hand or imported from
an iCalendar file that the college does not publish. They have no notion of roles, batches,
rooms or room availability.

\paragraph{College ERP / LMS timetable modules.}
ERP products and learning-management systems often include a timetable module, but the data must be
keyed in by staff, the modules are usually closed and licensed per student, and they rarely offer
free-room search, clash reports or natural-language queries.

\paragraph{Natural-language interfaces built on large language models.}
Recent assistants translate questions into SQL with large language models (text-to-SQL). They are
flexible, but answers are not reproducible, they can produce plausible but wrong queries, and they
are vulnerable to prompt-injection attacks that lead to SQL injection~\cite{pedro2023,greshake2023}.

\subsection{Review of Research Papers and Technical Resources}
Table~\ref{tab:litreview} summarises the papers and technical resources reviewed, grouped by
theme. The rest of this section discusses the most important ones.

\paragraph{Timetabling.}
Schaerf's survey~\cite{schaerf1999} classifies school, course and examination timetabling problems
and the search techniques used to solve them. Burke and Petrovic~\cite{burke2002} review research
directions in university timetabling, including constraint-based, meta-heuristic and hybrid methods,
and point out the gap between research prototypes and real institutional data. Lewis~\cite{lewis2008}
surveys meta-heuristics for university course timetabling and separates hard constraints (no clashes
of teachers, rooms or student groups) from soft ones. These works define the clash model we
\emph{check} after publication (professor, room and batch overlap), but they all assume that the
timetable is created by the software. None of them addresses reading an already published timetable.

\paragraph{Table recognition and OCR.}
Zanibbi, Blostein and Cordy~\cite{zanibbi2004} survey table recognition and describe the pipeline of
table detection, structure recognition (rows, columns, spanning cells) and interpretation. They note
that interpretation is domain specific. Kieninger and Dengel~\cite{kieninger1998} (T-Recs) recognise
tables from word bounding boxes without ruling lines. We use the same idea in our geometric
fallback, which rebuilds columns from whitespace gutters. The ICDAR 2013 table competition~\cite{gobel2013}
established standard evaluation of table location and structure, showing that cell-structure recovery
remains error-prone and motivating our per-cell verification status. Smith~\cite{smith2007} describes
the Tesseract OCR engine (adaptive classifier, line finding and word recognition) used for scanned
pages and photographs. PyMuPDF~\cite{pymupdf} provides text with coordinates and ruled-table detection
that preserves merged-cell spans.

\paragraph{Natural-language interfaces to databases.}
Androutsopoulos, Ritchie and Thanisch~\cite{androutsopoulos1995} review classical natural-language
interfaces to databases (NLIDB), including pattern-matching, syntax-based and semantic-grammar
systems. They note that a narrow domain allows simple, reliable pattern-based interpretation. Seq2SQL~\cite{zhong2017}
and the Spider benchmark~\cite{yu2018} show that neural text-to-SQL can generalise across schemas, but
exact-match accuracy on complex queries remains well below 100\%, which is unacceptable when a student
relies on the answer to reach a lab on time. Greshake et~al.~\cite{greshake2023} and Pedro
et~al.~\cite{pedro2023} demonstrate prompt-injection attacks on LLM-integrated applications, including
attacks that turn prompts into SQL injection. These findings led us to a closed intent registry with
parameterised templates instead of generated SQL.

\paragraph{Time intervals and conflicts.}
Allen's interval algebra~\cite{allen1983} defines the thirteen relations between time intervals. Our
clash rule uses the \emph{overlaps/during/equals} family, $a.\mathit{start} < b.\mathit{end} \wedge
b.\mathit{start} < a.\mathit{end}$, with start-inclusive, end-exclusive slots, so consecutive lectures do
not clash.

\paragraph{Security and architecture.}
Sandhu et~al.~\cite{sandhu1996} define role-based access control (RBAC), which underlies our five
roles and their data scopes. JSON Web Tokens are specified in RFC~7519~\cite{rfc7519}. Argon2, the
winner of the Password Hashing Competition, is a memory-hard function resistant to GPU
cracking~\cite{biryukov2016}. The OWASP Top~10~\cite{owasp2021} lists broken access control and
injection among the most critical web risks, and it guided our IDOR and injection controls.
Fielding's REST dissertation~\cite{fielding2000} motivates the stateless resource-oriented API, and
Codd's relational model~\cite{codd1970} with normalisation theory~\cite{silberschatz2019} underlies
the schema. PostgreSQL's \texttt{SELECT \ldots{} FOR UPDATE SKIP LOCKED}~\cite{postgresdocs} lets
several worker processes take jobs from one table without a separate message broker.

\begin{longtable}{@{}p{0.6cm}p{3.6cm}p{5.4cm}p{5.6cm}@{}}
\caption{Summary of reviewed literature and resources}\label{tab:litreview}\\
\toprule
\textbf{\#} & \textbf{Work} & \textbf{Contribution} & \textbf{Relevance / limitation for our problem}\\
\midrule
\endfirsthead
\toprule
\textbf{\#} & \textbf{Work} & \textbf{Contribution} & \textbf{Relevance / limitation for our problem}\\
\midrule
\endhead
1 & Schaerf, 1999~\cite{schaerf1999} & Survey of school, course and exam timetabling formulations and algorithms & Defines clash constraints; assumes timetable is generated, not read\\
2 & Burke \& Petrovic, 2002~\cite{burke2002} & Research directions in university timetabling & Highlights gap between research and real institutional data\\
3 & Lewis, 2008~\cite{lewis2008} & Survey of meta-heuristics for university timetabling; hard vs soft constraints & Hard constraints become our conflict checks\\
4 & Zanibbi et al., 2004~\cite{zanibbi2004} & Survey of table recognition: detection, structure, interpretation & Interpretation is domain specific, so we built a timetable-specific grid parser\\
5 & Kieninger \& Dengel, 1998~\cite{kieninger1998} & T-Recs: tables from word boxes without ruling lines & Basis of our geometric (gutter-based) fallback\\
6 & G\"obel et al., 2013~\cite{gobel2013} & ICDAR 2013 table competition and metrics & Structure recovery is error-prone, so we assign a status to each entry\\
7 & Smith, 2007~\cite{smith2007} & Architecture of the Tesseract OCR engine & OCR for scans/photos; OCR entries capped at \textsc{Unverified}\\
8 & Androutsopoulos et al., 1995~\cite{androutsopoulos1995} & Review of NL interfaces to databases & Narrow domain allows reliable rule-based interpretation\\
9 & Zhong et al., 2017~\cite{zhong2017} & Seq2SQL: neural text-to-SQL with RL & Flexible but not exact; needs training data\\
10 & Yu et al., 2018~\cite{yu2018} & Spider cross-domain text-to-SQL benchmark & Accuracy on complex queries far below 100\%\\
11 & Greshake et al., 2023~\cite{greshake2023} & Indirect prompt injection in LLM-integrated apps & Risk of letting an LLM drive data access\\
12 & Pedro et al., 2023~\cite{pedro2023} & Prompt-to-SQL injection in LLM web apps & Motivates templates instead of generated SQL\\
13 & Allen, 1983~\cite{allen1983} & Interval algebra for temporal reasoning & Formal basis of the overlap/clash rule\\
14 & Sandhu et al., 1996~\cite{sandhu1996} & Role-based access control models & Five roles with data scopes\\
15 & Biryukov et al., 2016~\cite{biryukov2016} & Argon2 memory-hard password hashing & Password storage\\
16 & Fielding, 2000~\cite{fielding2000} & REST architectural style & Stateless JSON API design\\
17 & Codd, 1970~\cite{codd1970} & Relational model of data & Normalised relational schema\\
18 & RFC 7519~\cite{rfc7519}, OWASP Top 10~\cite{owasp2021} & JWT standard; top web security risks & Token format; IDOR and injection controls\\
19 & FET, UniTime, aSc~\cite{fet,unitime,asc} & Timetable generators & Construct timetables; no reading, search or role views\\
20 & Tabula, Camelot, PyMuPDF~\cite{tabula,camelot,pymupdf} & PDF table extraction tools & Raw cells only; no timetable semantics\\
\bottomrule
\end{longtable}

\subsection{Comparative Analysis}
Table~\ref{tab:compare} compares the existing approaches on criteria that matter to the users of a
published college timetable.

\begin{table}[H]
\centering
\caption{Comparison of existing approaches with the proposed system}
\label{tab:compare}
\footnotesize
\setlength{\tabcolsep}{3pt}
\begin{tabular}{@{}p{3.6cm}*{7}{>{\centering\arraybackslash}p{1.6cm}}@{}}
\toprule
\textbf{Criterion} & \textbf{Manual PDF} & \textbf{FET / UniTime / aSc} & \textbf{Tabula / Camelot} & \textbf{Calendar apps} & \textbf{ERP module} & \textbf{LLM text-to-SQL} & \textbf{Proposed}\\
\midrule
Reads published PDF/photo & -- & No & Partly & No & No & No & \textbf{Yes}\\
Timetable semantics (labs, legends) & Human & Yes & No & No & Yes & Depends & \textbf{Yes}\\
Role-based views & No & Limited & No & No & Yes & No & \textbf{Yes (5)}\\
Plain-English questions & No & No & No & Limited & No & Yes & \textbf{Yes (composed plans)}\\
Reproducible, injection-safe answers & -- & -- & -- & -- & Yes & No & \textbf{Yes}\\
Clash detection & No & At creation & No & No & Sometimes & No & \textbf{Yes, after publication}\\
Version archive & Ad hoc & Files & No & No & Sometimes & No & \textbf{Yes, audited}\\
Personal timetable kept private & -- & No & No & Yes & No & No & \textbf{Yes}\\
Cost & Free & Free/Paid & Free & Free & Paid & Paid API & \textbf{Free (open source)}\\
Mobile use & PDF viewer & Limited & No & Yes & Varies & Varies & \textbf{Android, iOS, web}\\
\bottomrule
\end{tabular}
\end{table}

\subsection{Limitations of Existing Systems}
\begin{enumerate}
\item Generators (FET, UniTime, aSc) stop at a printable grid; the people who follow the timetable
      cannot search or personalise it.
\item Generic table extractors do not understand timetable conventions: stacked labs become one
      garbled string, legends are not applied, and times stay ambiguous.
\item Calendars and ERP modules require manual data entry, which is error-prone and is rarely
      repeated when the timetable is revised.
\item No existing tool used by our college detects clashes in the \emph{published} timetable.
\item LLM-based interfaces are not reproducible and introduce injection risks~\cite{pedro2023}.
\end{enumerate}

\subsection{Identified Engineering Gap}
There is no free system that takes the timetable \emph{as the college actually publishes it} and
turns it into trustworthy, role-scoped, searchable data with provenance, clash checks and version
history, while answering everyday questions in plain English safely. The gap is not a better
timetable-construction algorithm. It is the missing link between the published document and its
daily users.

\subsection{Proposed Improvements}
\begin{itemize}
\item A timetable-aware extraction pipeline with geometry, OCR and office adapters feeding one grid
      model, with per-entry provenance and verification status (closes limitations 2 and 3).
\item Admin review, correction and activation with an audited archive (closes limitation 3).
\item Interval-based clash detection that understands parallel lab groups and elective options
      (closes limitation 4).
\item A compositional NL query planner (typed plans compiled into parameterised queries) with clarification instead of guessing
      (closes limitation 5).
\item Role dashboards and a single code base for Android, iOS and the web (closes limitation 1).
\end{itemize}

\newpage
